%-----------------------Homework------------------------------------
%-------------------Arman Shokrollahi---------------------------------
%---------------------Coding Theory-------------------------------

\documentclass[a4 paper]{article}
% Set target color model to RGB
\input{preamble.tex}

\begin{document}
\homework{Práctia \#5}{Mayo. 13 al 14 de 2019}{MSc. Fausto M. Lagos S.}{}{Aquí va su nombre}{}

En esta práctica se desarrollarán tres planteamientos cada uno con una valoración de dos puntos. El objetivo es afianzar sus conocimientos antes de la Prueba de Evaluación Continua correspondiente, puede utilizar todo el material bibliográfica a su disposición y también preguntar todo lo que considere necesario. Preste atención al margen derecho donde encontrará premios y bonificaciones adicionales.

\vspace{5mm}

\problem{1} \textbf{2 puntos} - Del Cálculo en una variable de George Thomas, lea la sección Sección 2.5 (pág 115) y resuelva el ejercicio 36 y 48
\vskip 3mm
\solution{} 
\vspace{30mm}

\problem{2} \textbf{2 puntos} - Del Cálculo en una variable de George Thomas, pág. 146 lea los ejemplos en \textbf{Generalización de límites que involucran $\dfrac{\sin \theta}{\theta}$} y resuelva los ejercicios 29 y 30.
\vskip 3mm
\solution{}
\vspace{30mm}

\problem{3} \textbf{2 puntos} - Del Cálculo de una variable de James Stewart resuelva el ejercicio 42 de la página 129.
\vskip 3mm
\solution{}
\vspace{30mm}

\end{document}
